\chapter{LITERATURE REVIEW}

\section{E-Commerce Competitive Intelligence}

Competitive intelligence (CI) in e-commerce refers to the systematic collection and
analysis of publicly available market data to inform strategic decisions \cite{gilad1986}.
Unlike traditional industries where CI relies on surveys or industry reports, digital
marketplaces generate continuous, machine-readable signals—pricing data, rankings,
reviews, and inventory status—that enable near-real-time competitive monitoring.

Early work by \cite{chen2010} demonstrated that automated price monitoring on B2C
platforms could detect competitor repricing events with sub-hour latency, enabling
dynamic pricing responses. More recently, \cite{liu2023marketplace} highlighted that
Amazon's A9 ranking algorithm creates feedback loops where Best Seller Rank (BSR)
itself influences organic traffic, making BSR monitoring a leading indicator of
competitive threat.

\section{Amazon Marketplace Dynamics}

Amazon's marketplace exhibits several characteristics that distinguish it from
traditional retail environments:

\textbf{Best Seller Rank (BSR)} is Amazon's proprietary sales velocity indicator,
updated hourly and calculated separately within each product category
\cite{amazon2024bsr}. Empirical research by \cite{chevalier2006} found a strong
negative correlation between BSR and sales volume, though the exact formula remains
undisclosed. BSR has been validated as a proxy for relative sales performance across
multiple studies \cite{sopadjieva2017}.

\textbf{Buy Box} ownership determines which seller's offer is displayed as the primary
purchase option. \cite{chen2016buybox} found that the Buy Box winner captures
approximately 82\% of sales volume for a given ASIN, making Buy Box monitoring a
critical competitive signal.

\textbf{Listing quality} has been linked to conversion rates by \cite{ghose2009},
who found that richer product content (more images, longer descriptions, higher review
counts) correlated with higher purchase probability. Amazon's own A+ Content program
(also known as Enhanced Brand Content, EBC) was introduced to formalize these
content quality standards.

\section{Web Scraping for Market Research}

Web scraping has emerged as a primary data collection method in e-commerce research.
\cite{lewis2011} demonstrated the use of automated scrapers for price comparison
across retail websites. More recent work has leveraged cloud-based scraping
infrastructure to overcome IP blocking and CAPTCHA challenges common in large-scale
data collection \cite{mitchell2018webscraping}.

For Amazon specifically, scraping approaches must navigate anti-bot mechanisms including
request throttling, CAPTCHA injection, and session fingerprinting. Solutions such as
residential proxy rotation and headless browser simulation have been employed in
academic contexts \cite{richardson2019scrapy}.

In this thesis, we utilize the Apify cloud platform with the junglee/Amazon-crawler
actor, which provides managed proxy rotation and browser simulation, abstracting the
infrastructure complexity of direct scraping.

\section{Perceptual Hashing for Image Change Detection}

Perceptual hashing (pHash) algorithms generate compact fingerprints of images that
are robust to minor transformations (compression, resizing) while remaining sensitive
to meaningful content changes \cite{zauner2010phash}. The Hamming distance between
two pHash values provides a measure of visual similarity, with distances above a
threshold (typically 10–15 bits for a 64-bit hash) indicating substantive image
changes.

In the context of Amazon listings, product image changes can signal listing hijacking
(unauthorized sellers replacing images), A/B testing of hero images, or seasonal
packaging updates. To our knowledge, this thesis presents one of the first applications
of pHash-based monitoring to Amazon listing change detection.

\section{Sentiment Analysis of Product Reviews}

Customer review sentiment analysis has been extensively studied in the e-commerce
context. \cite{liu2012sentiment} provided a foundational taxonomy of opinion mining
approaches, distinguishing document-level, sentence-level, and aspect-level sentiment
classification.

Early approaches used lexicon-based methods (e.g., VADER \cite{hutto2014vader}),
which assign polarity scores based on predefined word lists. While computationally
efficient, these methods struggle with domain-specific language and sarcasm.

Transformer-based models, introduced by \cite{vaswani2017attention}, revolutionized
NLP by enabling contextual language understanding. BERT \cite{devlin2019bert} and its
variants have achieved state-of-the-art performance on sentiment classification tasks.
In particular, RoBERTa \cite{liu2019roberta} — a robustly optimized version of BERT
trained with larger batches and more data — has demonstrated superior performance on
social media and review sentiment tasks.

The \texttt{cardiffnlp/twitter-roberta-base-sentiment-latest} model, fine-tuned on
Twitter data with three-class classification (positive, neutral, negative), has been
validated on multiple review domains \cite{loureiro2022timelms} and serves as the
backbone of our sentiment pipeline.

\section{Price Segmentation with Clustering}

K-Means clustering \cite{macqueen1967kmeans} partitions observations into $k$ groups
by minimizing within-cluster variance. In retail analytics, price-based K-Means
clustering has been applied to market segmentation \cite{wedel2000marketSegmentation},
identifying natural price tiers that correspond to consumer value perceptions.

For Amazon product categories, price tiers are rarely explicit but emerge from
competitive dynamics. \cite{smith2019pricing} found that consumer electronics
categories typically exhibit bimodal or trimodal price distributions corresponding
to entry-level, mid-range, and premium segments, making $k=3$ a natural choice.

\section{Time Series Forecasting}

Price forecasting in e-commerce has been approached through both classical econometric
and machine learning methods. ARIMA models \cite{box2015time} provide interpretable
forecasts but require stationarity and substantial historical data (typically $n > 50$).

Facebook Prophet \cite{taylor2018forecasting} addresses the short-series limitation
through a decomposable additive model incorporating trend, seasonality, and holiday
effects. Prophet has been shown to perform competitively with ARIMA on datasets with
as few as 30 observations \cite{triebe2021neuralprophet}, making it appropriate for
our 15-day data collection window.

\section{Research Gap}

The existing literature has addressed individual components of Amazon marketplace
analysis — BSR dynamics, review sentiment, price segmentation — in isolation. However,
no prior work has integrated these analytical streams into a unified, automated,
end-to-end pipeline with daily refresh cadence and real-time alerting. This thesis
contributes such an integration, providing a replicable framework for practitioners
and researchers in the e-commerce competitive intelligence domain.
